% Expose IX, sections 8.2--8.3.
% French transcription lines 2159--2188; authority scan physical pp. 428--429
% (scan indices 427--428; source folios 416--417). Both pages were checked
% directly in four overlapping 1100-dpi bands; physical p. 430 was checked
% for the exact section-9 boundary. The authority's evident A_K -> A'_K
% polarization is restored where the working transcription misplaces the
% prime as A_{K'}; its evident "extensions ... de K'" slip is rendered
% according to the mathematically required extensions of K.

\medskip
\noindent
8.2. The hypotheses are those of 8.1. It is clear that the extension
(8.1.1) depends functorially on the abelian scheme $A_K$, and that
its formation commutes with every morphism of traits
$\widetilde S\longrightarrow S$. Choose a polarization of $A_K$,
hence an isogeny $A_K\longrightarrow A'_K$. It induces a
homomorphism from (8.1.1) to (8.1.6), which is immediately seen, by
looking at the $T_\ell$, to be an isogeny term by term. Likewise,
using 7.5 and descent, one sees that every biextension by
$\mathbb{G}_{mK}$ of abelian schemes $A_K$ and $\overline A_K$ over
$K$ defines a biextension of $(B_K,\overline B_K)$ by
$\mathbb{G}_{mK}$, where $B_K$ and $\overline B_K$ denote the
abelian quotients of the Raynaud extensions corresponding to $A_K$
and $\overline A_K$, respectively. When $W$ is a perfect duality, so
is the corresponding biextension of $(B_K,\overline B_K)$. In
particular, with the notation of 7.6, the abelian schemes $B_K$ and
$B'_K$ are canonically dual to one another.

The correspondence between biextensions on $(A_K,\overline A_K)$
and biextensions on $(B_K,\overline B_K)$ has evident functoriality
properties as $A_K$ and $\overline A_K$ vary. These follow by descent
from the analogous properties stated in 7.5, and we leave their
formulation to the reader. We also leave to the reader the
compatibility of the preceding construction with every change of
traits $\widetilde S\longrightarrow S$.

\medskip
\noindent
8.3. Return to the group $A_{K,K'}^{\natural}$ and study its
variation with $K'$. It is clear that two isomorphic extensions
$K'$ and $K'_1$ of $K$ give canonically isomorphic groups, so it is
enough to consider the subextensions $K'$ of the separable closure
$\overline K$ of $K$. If $K'\hookrightarrow K''$ are two such finite
Galois subextensions of $\overline K$, it follows immediately from
3.3 that there is a canonical \underline{open immersion}

\begin{equation}\tag{8.3.1}
A^{\natural}_{K,K'} \hookrightarrow A^{\natural}_{K,K''}
 \qquad ,
\end{equation}
which identifies the first group with an open subgroup of the
second. As $K'$ varies, the $A^{\natural}_{K,K'}$ thus form an
inductive system of smooth groups of finite type over $K$, whose
transition morphisms are open immersions. If we set

\begin{equation}\tag{8.3.2}
\Phi_{K,K'}=A^{\natural}_{K,K'}
 /{A^{\natural}_{K,K'}}^{o} \qquad ,
\end{equation}
the $\Phi_{K,K'}$ themselves form an inductive system of finite
etale groups over $K$, which will be studied in greater detail
below.

We shall set

\begin{equation}\tag{8.3.3}
A^{\flat}_{K}=\varinjlim_{K'} A^{\natural}_{K,K'} \qquad ,
\end{equation}
where the limit is taken in the category of fppf sheaves over $K$.
Since the transition morphisms in this inductive system are open
immersions, it is clear moreover that this limit is represented by
a group scheme over $K$ (in general non-quasi-compact), which we
shall naturally denote by the same symbol $A^{\flat}_{K}$. For this
group scheme one has

\begin{equation}\tag{8.3.4}
{A_{K}^{\flat}}^{o}=A_{K}^{\natural o} \quad , \qquad
A^{\flat}_{K}/{A_{K}^{\flat}}^{o}
 \simeq \varinjlim_{K'}\Phi_{K,K'} \qquad .
\end{equation}
The latter group will be determined in 11.10: it is canonically
isomorphic to $\underline{M}_{K}\otimes\mathbb{Q}/\mathbb{Z}$, with
the notation of (8.1.8).
